\begin{ouverture}[calcul]{1}{Calculer avec les nombres relatifs}{Partie 1 · Nombres et calculs}
{Signes · Produits · Quotients · Priorités}
\objectif{Multiplier et diviser des nombres relatifs.}
\objectif{Organiser un calcul et contrôler son signe.}
\objectif{Interpréter un résultat dans une situation concrète.}
\prerequis{Addition, soustraction et nombres décimaux.}
\end{ouverture}
\pagestyle{otMain}\essentiel\fontsize{11}{15}\selectfont
\phantomsection\label{cours-1-1}
\essSection{1}{Comprendre les signes}
\begin{definition}L’opposé de $a$ est $-a$ : leur somme vaut 0. La distance à zéro est la valeur du nombre sans son signe. Soustraire un nombre revient à ajouter son opposé : $a-b=a+(-b)$.\end{definition}
\begin{propriete}Pour multiplier deux nombres non nuls, multiplier leurs distances à zéro. Le résultat est positif si leurs signes sont identiques, négatif sinon. Un produit avec un facteur nul est nul. Un produit de plusieurs facteurs non nuls est négatif si le nombre de facteurs négatifs est impair.\end{propriete}
\begin{exemple}$(-4)\times(-2{,}5)=10$ ; $(-3)\times5=-15$. Dans $(-2)\times(-3)\times(-4)$, trois facteurs sont négatifs : le produit vaut $-24$.\end{exemple}
\begin{attention}Deux signes « moins » n’ont pas toujours le même rôle. Dans $-4-3$, on soustrait 3 : le résultat est $-7$. Dans $(-4)\times(-3)$, on multiplie deux nombres négatifs : le résultat est 12.\end{attention}
\clearpage
\phantomsection\label{cours-1-2}
\essSection{2}{Diviser et organiser}
\begin{propriete}Le quotient $a/b$ est le nombre qui, multiplié par $b$, donne $a$, avec $b\ne0$. Sa règle de signe est celle du produit. On ne divise jamais par zéro.\end{propriete}
\begin{savoirfaire}[Une ligne par étape]Calculer les parenthèses, puis les puissances, puis produits et quotients de gauche à droite, enfin sommes et différences de gauche à droite. Réécrire les parties inchangées pour conserver la structure.\end{savoirfaire}
\begin{exemple}$6-2\times(-3+7)=6-2\times4=6-8=-2$. À l’inverse, $(6-2)\times(-3+7)=4\times4=16$ : les parenthèses changent le calcul.\end{exemple}
\begin{savoirfaire}[Contrôler]Prévoir le signe d’un produit, estimer l’ordre de grandeur, puis refaire une opération inverse. Pour un problème, préciser le sens du signe et l’unité : $-8$ euros représente une diminution de 8 euros.\end{savoirfaire}

\clearpage\rubrique{Découverte}
\begin{blocLivre}{ex-1-1}\begin{exercice}[1]{Un produit positif}
Calculer $(-7)\times(-6)$.
\end{exercice}\end{blocLivre}
\begin{blocLivre}{ex-1-2}\begin{exercice}[1]{Un produit négatif}
Calculer $2{,}4\times(-5)$.
\end{exercice}\end{blocLivre}
\begin{blocLivre}{ex-1-3}\begin{exercice}[1]{Un quotient}
Calculer $(-54):(-9)$ et $18:(-4)$.
\end{exercice}\end{blocLivre}
\begin{blocLivre}{ex-1-4}\begin{exercice}[1]{Opposé ou produit}
Calculer $-8-(-3)$ et $(-8)\times(-3)$.
\end{exercice}\end{blocLivre}
\begin{blocLivre}{ex-1-5}\begin{exercice}[1]{Zéro et division}
Calculer $0:(-12)$. Peut-on calculer $(-12):0$ ?
\end{exercice}\end{blocLivre}

\rubrique{Entraînement}
\begin{blocLivre}{ex-1-6}\begin{exercice}[2]{Compter les signes}
Sans effectuer les multiplications, donner le signe de $(-2)\times5\times(-3)\times(-7)$, puis calculer.
\end{exercice}\end{blocLivre}
\begin{blocLivre}{ex-1-7}\begin{exercice}[2]{Respecter les priorités}
Calculer $A=14-3\times(-2)$ et $B=(14-3)\times(-2)$.
\end{exercice}\end{blocLivre}
\begin{blocLivre}{ex-1-8}\begin{exercice}[2]{Une expression imbriquée}
Calculer $18:[-2\times(5-8)]$.
\end{exercice}\end{blocLivre}
\begin{blocLivre}{ex-1-9}\begin{exercice}[2]{Calculer efficacement}
Calculer $(-25)\times17\times(-4)$ en regroupant les facteurs.
\end{exercice}\end{blocLivre}
\begin{blocLivre}{ex-1-10}\begin{exercice}[2]{Une erreur de calcul}
Nora écrit $24:(-6)\times2=24:(-12)=-2$. Corriger.
\end{exercice}\end{blocLivre}

\rubrique{Problèmes}
\begin{blocLivre}{ex-1-11}\begin{exercice}[3]{Évolution de température}
À midi, il fait $3\,^\circ$C. Pendant 5 heures, la température varie de $-1{,}6\,^\circ$C par heure. Quelle température obtient-on ?
\end{exercice}\end{blocLivre}
\begin{blocLivre}{ex-1-12}\begin{exercice}[3]{Un compte de jeu}
On part de 20 points. Trois erreurs font perdre chacune 7 points et deux réussites rapportent chacune 4 points. Donner le bilan final.
\end{exercice}\end{blocLivre}
\begin{blocLivre}{ex-1-13}\begin{exercice}[3]{Un nombre inconnu}
Quel nombre faut-il multiplier par $-2{,}5$ pour obtenir 15 ? Vérifier.
\end{exercice}\end{blocLivre}
\begin{blocLivre}{ex-1-14}\begin{exercice}[3]{La machine à signes}
Le produit de 19 nombres négatifs non nuls est-il positif ? Que devient son signe si l’on remplace un facteur par son opposé, puis si l’on remplace un facteur par zéro ?
\end{exercice}\end{blocLivre}
\begin{blocLivre}{ex-1-15}\begin{exercice}[3]{Un calcul à contrôler}
Calculer $C=-3[4-2(-5+8)]+12:(-4)$. Donner un contrôle du signe final.
\end{exercice}\end{blocLivre}

\clearpage
